\documentclass{article}
\usepackage{../../third_party/iclr2027/iclr2027_conference,times}
\input{../../third_party/iclr2027/math_commands.tex}
% Generated by scripts/build_paper_claim_summary.py; do not edit by hand.
\newcommand{\PTLPrimaryDeltaCross}{0.324175}
\newcommand{\PTLPrimaryDeltaAdj}{0.093331}
\newcommand{\PTLPrimaryDeltaAdjCI}{[0.056748, 0.134389]}
\newcommand{\PTLPrimarySystemaCross}{0.110196}
\newcommand{\PTLPrimarySystemaAdj}{0.035181}
\newcommand{\PTLPrimarySystemaAdjCI}{[0.017110, 0.057823]}
\newcommand{\PTLPrimaryRankCross}{0.328296}
\newcommand{\PTLPrimaryRankAdj}{0.073300}
\newcommand{\PTLPrimaryRankAdjCI}{[0.035638, 0.111859]}
\newcommand{\PTLDecisionRho}{0.775026}
\newcommand{\PTLDecisionRhoDisplay}{0.78}
\newcommand{\PTLDecisionRhoCI}{[0.73, 0.89]}
\newcommand{\PTLWithinMetricRho}{0.26}
\newcommand{\PTLWithinMetricRhoCI}{[0.02, 0.45]}
\newcommand{\PTLMeanFidelityShiftRhoD}{0.69}
\newcommand{\PTLMeanFidelityShiftRhoRegret}{0.77}
\newcommand{\PTLSameContextRange}{0.0012--0.0054}
\newcommand{\PTLNadigDelta}{0.176}
\newcommand{\PTLNadigSystema}{0.138}
\newcommand{\PTLNadigRank}{0.129}
\newcommand{\PTLFrangiehLabels}{243}

% Generated by scripts/build_scientific_upgrade_tables.py; do not edit by hand.
\newcommand{\PTLDecompositionIdentityError}{8.33e-17}
\newcommand{\PTLDecompositionVsCanonicalMaxDiff}{0.022}
\newcommand{\PTLComparatorKendallRho}{0.87}
\newcommand{\PTLComparatorJaccardRho}{0.84}
\newcommand{\PTLIncrementalMeanDeltaMAE}{0.013}
\newcommand{\PTLIncrementalMeanDeltaMAECI}{[-0.001, 0.027]}
\newcommand{\PTLIncrementalDAdjMetricOnlyDeltaMAE}{-0.005}
\newcommand{\PTLIncrementalDAdjMetricOnlyDeltaMAECI}{[-0.019, 0.008]}
\newcommand{\PTLIncrementalDAdjBeyondMeanDeltaMAE}{-0.005}
\newcommand{\PTLIncrementalDAdjBeyondMeanDeltaMAECI}{[-0.014, 0.003]}
\newcommand{\PTLTargetAuditRows}{360}
\newcommand{\PTLTargetAuditFalseReuse}{0.019}

% Generated by scripts/write_v3_manuscript_tables.py; do not edit by hand.
\newcommand{\PTLNullDAdjFTR}{0.017}
\newcommand{\PTLNullKendallFTR}{0.952}
\newcommand{\PTLNullSpearmanFTR}{0.995}
\newcommand{\PTLNullKendallMatchedFTR}{0.003}
\newcommand{\PTLNullJaccardFTR}{1.000}
\newcommand{\PTLNullJaccardMatchedFTR}{0.000}
\newcommand{\PTLAltDAdjPower}{1.000}
\newcommand{\PTLAltKendallMatchedPower}{0.997}
\newcommand{\PTLAltJaccardMatchedPower}{0.528}

\usepackage{graphicx}
\usepackage{hyperref}
\usepackage{url}
\usepackage{booktabs}
\usepackage{placeins}
\hypersetup{hidelinks}

\title{Finite-Measurement Transport of Perturbation Reliability Rankings}
\author{Anonymous Authors}

\begin{document}
\maketitle

\begin{abstract}
Predictions validated in one biological context are often considered for use
in another, yet perturbations with low prediction error in one context need not
retain their relative accuracy after a context change. We study this ordering
of perturbation-level realized error while the predictor is held fixed. Because
finite cell counts can themselves induce apparent rank changes,
we define a tie-aware discrepancy that compares pairwise ordering
distributions and subtracts expected within-context ordering disagreement. The
resulting finite-measurement estimand is zero when the two pair-state laws
agree, while its unbiased estimator may be negative at finite sample size. In the
Frangieh Perturb-CITE-seq dataset \citep{frangieh2021multimodal}, reliability
rankings change across control, co-culture and IFN$\gamma$ conditions under
three complementary risk metrics. The changes exceed same-context sampling
baselines, become increasingly resolvable with cell budget, and are
reproduced in independent HepG2--Jurkat and iPSC--neuron raw-cell analyses
\citep{nadig2025trade,tian2019crispri}. Matched source-only RBF, direct-MLP, latent-MLP and
strict source-frozen GEARS sensitivities preserve positive transport signs on
their respective eligible Frangieh surfaces, although their within-metric profiles are not
identical. Reliability reordering also changes finite top-$k$ experimental
shortlists. In the tested controls, however, $D_{\rm adj}$ did not improve
held-out regret prediction beyond metric identity and mean-risk shift. A known-truth
comparator audit shows why the finite-measurement correction matters: uncorrected
rank statistics can declare transport under a same-context null, whereas
matched-null corrections restore calibration. The framework therefore supports
measurement-aware auditing of cross-context changes in model-error ordering.
\end{abstract}

\section{Introduction}
Single-cell perturbation screens make it possible to measure many interventions
in parallel, but exhaustively measuring them across biological contexts is
often impractical. Pooled Perturb-seq established a route from genetic
perturbations to high-dimensional single-cell readouts
\citep{dixit2016perturbseq,adamson2016parseq}; genome-scale screens and
harmonized perturbation resources have since expanded both the intervention
universe and the range of measurement depths \citep{replogle2022genome,
peidli2024scperturb}. A computational predictor can therefore help determine
which source-validated predictions may be reused after a context change.

That narrower reuse decision is an ordering problem. Existing perturbation-prediction studies
usually ask whether a model reconstructs an unseen response or generalizes to
an unseen perturbation. Representative approaches span latent-variable,
optimal-transport, graph-based and foundation-model methods, including scGen,
CPA, CellOT, GEARS, scGPT and TxPert
\citep{lotfollahi2019scgen,lotfollahi2023cpa,bunne2023cellot,
roohani2023gears,cui2024scgpt,wenkel2026txpert}. Benchmarks broaden this
comparison across datasets and metrics \citep{generalizable2025benchmark},
while direct audits show that complex models need not outperform simple linear
baselines under a fixed protocol \citep{ahlmann2025simple,wong2025simple}.
These studies evaluate response fidelity; they do not directly ask whether
the identity of reliably predicted perturbations survives a biological context
shift.

A ranking shift is not automatically a biological-context effect. With finite
numbers of cells, perturbation risks fluctuate, and two observed orders can
differ even when the underlying order is not resolved. The number of cells per
perturbation--context pseudoreplicate is therefore part of the inferential problem, not merely
technical details \citep{peidli2024scperturb,squair2021augur}. Work on ranking
under noisy measurements makes the same point for uncertain and tied orders
\citep{zuk2007ranking,sankaran2025ranking}. Our question is complementary to
calibration and selective prediction: after accounting for this finite-cell
ordering floor, does the relative ordering of realized model risk change
between biological contexts?

We answer this question with a source-frozen protocol. Predictions are fixed
before target outcomes are observed; context-specific risks are then compared
through a tie-aware, measurement-adjusted discrepancy. For each perturbation
pair, the estimator compares the distributions of lower, tied and higher risk
states across contexts and subtracts expected within-context ordering disagreement
using an unbiased order-2 U-statistic~\citep{hoeffding1948class}. The primary
analysis uses the Frangieh context surface, same-context pseudoreplicate
controls, independent HepG2--Jurkat and iPSC--neuron raw-cell replications,
matched-cell-budget resampling and a finite top-$k$ decision layer.

\textbf{We make three contributions.} First, we formalize reliability transport
as a finite-measurement comparison of perturbation-pair ordering distributions
rather than mean response fidelity. Second, we show that context-associated
reordering exceeds a matched sampling floor, has an explicit measurement
boundary, and recurs across independent biological systems and matched
nonlinear-predictor sensitivities. Third, we quantify its consequence for
finite experimental shortlists while keeping post-hoc transport measurement
separate from prospective source-only prediction.

\begin{figure*}[t]
  \centering
  \includegraphics[width=\textwidth]{../../results/figures/reliability_transportability/reliability_transportability_fig1_graphical_abstract.pdf}
  \caption{\textbf{Reliability ordering under biological context shift.}
  The source predictor is frozen before evaluation in the target context.
  The rank braid retains all 243 background rank trajectories and highlights
  a deterministic set of source-rank quantile exemplars for legibility. The
  compact shortlist preview reports 25 perturbations per context:
  6 retained, 19 dropped and 19 newly entering the target shortlist. The full
  243-label accounting remains in the Fig.~1 source table, while the detailed
  replacement flow is reserved for Fig.~4A. The decomposition uses the same
  source-frozen predictor and prespecified evaluation protocol on its declared eligible
  transfer set. Metric colors are fixed throughout the paper.}
  \label{fig:framework}
\end{figure*}

\section{Problem setting and measurement-adjusted ordering}
Let $X_{p,e}^{(r)}$ denote the raw cells sampled for perturbation $p$ in context
$e$ and Monte Carlo replicate $r$. Each cell is library-size normalized to
$10^4$, transformed by $\log(1+x)$, and aggregated into the declared gene-panel
response $y_{p,e}^{(r)}$ (a perturbation-versus-control delta for delta-based
metrics). The primary predictor is an out-of-fold map: for perturbation $p$,
$f_{s,-\mathrm{fold}(p)}$ is fit only to labels outside $p$'s fixed fold in the
source training context $s$, using a ten-component PCA representation, ridge
penalty $\alpha=0.1$, and three independently seeded members. It produces a
fixed held-out vector $\hat y_{p,s}$ before any target outcome is examined. For
metric $m$, the realized risk is
\begin{equation}
 R_{p,e}^{(m,r)}=g_m(\hat y_{p,s},y_{p,e}^{(r)}),
\end{equation}
where larger values mean worse prediction and the direction $s\to e$ names the
training context and evaluation context, respectively. We use one minus cosine similarity
(delta cosine), one minus Systema centroid accuracy~\citep{systema2025}, and one minus
absolute-effect rank agreement; complete metric definitions, comparator
universes, and tie tolerances are given in Appendix~B. These risks are evaluated after target
outcomes are observed; source-only uncertainty is kept separate from the
ordering estimand.

For a pair $a=(p,q)$, define the tie-aware pair state
\begin{equation}
 Z_{a,e}^{(m,r)}=\operatorname{sign}\!\left(R_{p,e}^{(m,r)}-R_{q,e}^{(m,r)}\right)
 \in\{-1,0,+1\}.
\end{equation}
The resampling distribution gives $\pi_{a,e}^{(m)}=[\Pr(Z=-1),\Pr(Z=0),
\Pr(Z=+1)]$. For two contexts $e_1,e_2$,
\begin{equation}
  D_{\rm cross}=\mathbb E_{a\sim w}[1-\pi_{a,e_1}^{\mathsf T}\pi_{a,e_2}],\qquad
  D_{\rm within,e}=\mathbb E_{a\sim w}[1-\|\pi_{a,e}\|_2^2].
\end{equation}
The measurement-adjusted discrepancy, equivalently one half of the squared
delta-kernel MMD
\citep{gretton2012kernel}, is
\begin{equation}
 D_{\rm adj}=D_{\rm cross}-
 \frac{D_{\rm within,e_1}+D_{\rm within,e_2}}{2}
 =\frac12\mathbb E_a\|\pi_{a,e_1}-\pi_{a,e_2}\|_2^2
 \label{eq:ordering_identity}
\end{equation}
Here $w$ is uniform over the declared unordered perturbation-pair universe and
the delta kernel $k(z,z')=\mathbb 1\{z=z'\}$ compares the three-valued pair
states. Thus $D_{\rm adj}$ is the squared separation between the two ordering
distributions after removing their within-context coincidence terms. It is a
transport quantity for relative reliability, rather than an average change in
prediction accuracy.
For interpretability, write $t_e=\pi_{a,e}(0)$ and
$b_e=\pi_{a,e}(+1)-\pi_{a,e}(-1)$. The pairwise squared separation has the
exact decomposition
\begin{equation}
 D_{\rm adj}(a)=D_{\rm tie}(a)+D_{\rm strength}(a)+D_{\rm flip}(a),\quad
 \begin{aligned}
 D_{\rm tie}(a)&=\tfrac34(t_{e_1}-t_{e_2})^2,\\
 D_{\rm strength}(a)&=\tfrac14(|b_{e_1}|-|b_{e_2}|)^2,\\
 D_{\rm flip}(a)&=|b_{e_1}b_{e_2}|\,\mathbb{1}\{b_{e_1}b_{e_2}<0\}.
 \end{aligned}
 \label{eq:ordering_decomposition}
\end{equation}
For finite measurement, the reported components use a cross-product of
independent 15-seed halves while retaining both measurement replicates.
The empirical report is therefore $D_{\rm tie}^{\rm CF}$ plus a
directional-bias component; the nonlinear strength/flip split is retained as
population algebra, not presented as an unbiased finite-sample decomposition.
The canonical per-seed U-corrected $D_{\rm adj}$ is reconstructed from the
merged member-level risks. Finite-measurement component calibration and the
comparison between cross-fit and canonical estimates are reported in the
Supplement; they are not additional population identities.
\paragraph{Estimand properties.}
Because each $\pi_{a,e}$ is a probability vector, $0\leq D_{\rm adj}\leq 1$.
The discrepancy is symmetric in the two contexts and equals zero if and only
if the pair-state distributions agree almost everywhere over the declared
pair universe. It is therefore a squared divergence; we do not claim a
triangle inequality and use ``discrepancy'' rather than ``distance''.
The bound applies to the population estimand. Its unbiased finite-sample
U-corrected estimate $\widehat D_{\rm adj}$ is signed and is not clipped to
$[0,1]$; a negative estimate can occur when the observed pair-state
discrepancy is smaller than the finite-replicate correction. In the
deterministic tie-free limit, each pair law collapses to a point mass and
  $D_{\rm adj}$ reduces exactly to the normalized Kendall inversion fraction,
  $(1-\tau)/2$. More generally, it equals cross-context pair disagreement minus
  one half of each within-context pair disagreement, a population quantity that
  deterministic Kendall correlation does not define under repeated noisy
  measurement. Its distinct role is finite-measurement transport, where pair
states are uncertain and the matched within-context component must be
estimated rather than ignored.
This distinction is not cosmetic. Consider one perturbation pair whose strict
ordering probabilities are $(0.9,0,0.1)$ in the source and $(0.6,0,0.4)$ in
  the target. Both contexts have the same modal strict order, so Kendall after
  collapsing each pair law to its modal order records no inversion, whereas
$D_{\rm adj}(a)=\tfrac12[(0.9-0.6)^2+(0.1-0.4)^2]=0.09$. The estimand therefore
  detects transport of ordering reliability even when the modal order does not
  flip. Conversely, an expression-space MMD can be positive when every risk
ordering is unchanged; it tests a different population object. Hence
$D_{\rm adj}$ is neither renamed Kendall nor a generic distribution-shift
score: it is the declared average separation of finite-measurement pair-state
laws.
For finite resampling sets, the within-context term uses the unbiased order-2
U-statistic~\citep{hoeffding1948class}. A strict direction is declared stable when the 95\% equal-tail
  Beta(1,1) posterior credible interval excludes $0.5$, with at least eight
  non-tie Monte Carlo resampling outcomes supporting the strict comparison.
  This threshold diagnoses state-estimation stability and is not a count of
  biological replicates; otherwise the pair
remains unresolved. This categorical state rule is separate from the 90\%
intervals reported for effect estimates.
At the no-transport boundary the U-statistic is degenerate, so an ordinary
percentile bootstrap is not used as a zero-effect test. Detectability is
instead assessed from the empirical variance of the 30 independent
seed-level U-corrected contributions using a Student-$t$ interval. A
complete-seed block bootstrap is retained as a sensitivity analysis. This
preserves the registered seed as the independent unit and does not resample
the two context arms independently.

The primary Frangieh surface contains 243 shared perturbation labels across
three contexts. Fig.~\ref{fig:framework}B selects source-rank quantile
exemplars only to keep the complete braid legible, whereas its background
ribbons and the Fig.~1 source table retain the full 243-label source-frozen
surface. Fig.~1D is a compact decision preview; the detailed shortlist flow
appears only in Fig.~4A. Aggregate transport estimates use their declared full eligible
tables. The predictor uses a ten-component PCA basis, ridge penalty
$\alpha=0.1$, one global five-fold perturbation split and three independent
model members. Matched raw-cell pseudoreplicates use 30 resampling seeds and a
minimum of 20 cells per perturbation--context. The independent Nadig analysis
uses the exact shared HepG2--Jurkat perturbation set and is not pooled with
Frangieh.

\subsection{Decision consequences}
For this descriptive decision analysis, $\bar R_{p,e}^{(m)}$ is the mean
full-depth risk over the frozen model members and measurement replicates for
metric $m$. We set $k=\max(1,\lfloor q|\mathcal P|\rceil)$ for budget fraction
$q$ and break exact boundary ties by the fixed perturbation identifier. For a
source context $e_s$ and selection budget $k$, let
\begin{equation}
 S_k(e_s)=\operatorname*{arg\,min}_{S:\,|S|=k}\sum_{p\in S}\bar R_{p,e_s}^{(m)}
 \label{eq:source_selection}
\end{equation}
be the source-ranked shortlist. Its target-context loss is
$L_{e_t}(S)=k^{-1}\sum_{p\in S}\bar R_{p,e_t}^{(m)}$, and normalized selection regret
compares $L_{e_t}(S_k(e_s))$ with the target-oracle shortlist $S_k(e_t)$.
Thus the decision layer measures the consequence of transporting a source
reliability order, rather than evaluating a second predictor.

For the decision layer, the target regret at budget $k$ is defined as
\begin{equation}
 \operatorname{Regret}_k(e_s\!\to e_t)=
 L_{e_t}\!\left(S_k(e_s)\right)-L_{e_t}\!\left(S_k(e_t)\right),
 \label{eq:decision_regret}
\end{equation}
where the second term is the target-oracle loss. Let
$S_k^{\rm worst}(e_t)$ be the deterministic $k$-item set with largest target
risk and let $L_{e_t}(\mathcal U)$ be the mean target risk over the eligible
universe. We report two explicitly distinct normalizations,
\begin{equation}
 \operatorname{Regret}^{\rm rand}_k=
 \frac{\operatorname{Regret}_k}{L_{e_t}(\mathcal U)-L_{e_t}(S_k(e_t))},\qquad
 \operatorname{Regret}^{\rm worst}_k=
 \frac{\operatorname{Regret}_k}{L_{e_t}(S_k^{\rm worst})-L_{e_t}(S_k(e_t))}.
 \label{eq:normalized_regret}
\end{equation}
The first preserves the random-reference scale used in the primary result;
the second is bounded by the deterministic worst shortlist. Rows with a
non-positive normalization denominator are undefined and excluded; no such row
occurred in the reported surface. Both target-oracle quantities are descriptive
benchmarks computed from the same target table, not prospective policies.

\section{Results}
\subsection{Reliability ordering changes across biological contexts}
Figure~\ref{fig:framework} introduces the object of interest: predictions are
held fixed while target biology changes, so movement in the perturbation order
cannot be attributed to refitting the predictor. The full Frangieh atlas
contains 243 shared labels; Fig.~\ref{fig:framework}B retains all 243
background trajectories and highlights a deterministic set of source-rank
quantile exemplars so that individual ribbons remain readable.
Preprocessing and replication eligibility are summarized in Appendix A; the
complete candidate registry is retained with the repository artifacts.
Fig.~\ref{fig:framework}D previews the fixed top-10\% shortlist replacement example; the complete
label-level risk accounting is retained in the figure source table. The same-context resampling
control gives $D_{\rm adj}=\PTLSameContextRange$ across its nine context--metric
cells, whereas the matched biological transfers produce larger
measurement-adjusted values.

\begin{figure*}[t]
  \centering
  \includegraphics[width=\textwidth]{../../results/figures/reliability_transportability/reliability_transportability_fig2_transport_atlas.pdf}
  \caption{\textbf{Context-dependent reliability transport.} (A) Six directed
  Frangieh transfers with three metric estimates and 90\% intervals. (B)
  Stable-pair inversion fraction: fraction of pairs stable in both contexts
  whose strict ordering reverses. (C) Same-context independent
  pseudoreplicate control, shown as point estimates because it is a floor
  diagnostic rather than a biological transfer. (D) Independent Nadig
  HepG2--Jurkat replication in both source directions. Metric color is fixed;
  in panel D, marker shape identifies the source direction and context is
  otherwise carried by position and labels.}
  \label{fig:atlas}
\end{figure*}

The full-size Ctrl-to-IFN$\gamma$ transfer illustrates metric dependence. For
delta cosine, $D_{\rm cross}=\PTLPrimaryDeltaCross$ and
$D_{\rm adj}=\PTLPrimaryDeltaAdj$ ($\PTLPrimaryDeltaAdjCI$); for Systema
centroid, the corresponding values are $\PTLPrimarySystemaCross$ and
$\PTLPrimarySystemaAdj$ ($\PTLPrimarySystemaAdjCI$); for absolute-effect rank
they are $\PTLPrimaryRankCross$ and $\PTLPrimaryRankAdj$
($\PTLPrimaryRankAdjCI$). These generated values are read from the same
  prespecified analysis table used by the figure builder. The conservative joint-floor
stress test is retained separately, because it is not an additive component of
the measurement-adjusted estimand. The atlas in Fig.~\ref{fig:atlas} retains
metric and transfer heterogeneity rather than averaging it away.
The finite-sample estimator and tie rule are given in
Appendix B.

\subsection{Calibrated inference distinguishes transport from finite measurement}
The estimator addresses a finite-measurement target, so a useful comparison is
not only against raw rank distances but also against their strongest simple
matched-null variants. In a registered 120-condition known-truth simulation,
raw Kendall, Spearman and top-10 Jaccard declared transport in
\PTLNullKendallFTR{}, \PTLNullSpearmanFTR{} and \PTLNullJaccardFTR{} of null
conditions, respectively. Label bootstrapping alone did not remove this
finite-sample effect. Subtracting a matched within-context null reduced the
Kendall false-transport rate to \PTLNullKendallMatchedFTR{} and the Jaccard
rate to \PTLNullJaccardMatchedFTR{}, while the U-corrected $D_{\rm adj}$ rate
was \PTLNullDAdjFTR{}. At non-zero inversion fractions, detection power was
\PTLAltDAdjPower{} for $D_{\rm adj}$ and \PTLAltKendallMatchedPower{} for
matched-null Kendall. Thus the result is not that $D_{\rm adj}$ universally
dominates every rank statistic: precise measurements make deterministic rank
comparators reasonable, whereas $D_{\rm adj}$ supplies one coherent
pair-state population estimand with an explicit finite-replicate correction.
For the interval itself, direct variance estimation over the 30 independent
seed-level U-corrected contributions yielded a null false-transport rate of
$0.040$ with $0.960$ coverage in 2,880 known-truth trials, and $0.015$ with
$0.985$ coverage in 3,000 Frangieh semi-synthetic null trials. These intervals
are mildly conservative and replace the anti-conservative hierarchical
percentile intervals in every detectability result.
The full calibration table, including coverage, bias and the Jaccard
comparison, is provided in the Supplement.

\subsection{Independent raw-cell systems reproduce ordering transport}
The independent Nadig replication provides a different biological system and
context pair and a larger shared perturbation universe. The adjusted discrepancy
$D_{\rm adj}$ is \PTLNadigDelta{} for
delta cosine, \PTLNadigSystema{} for Systema centroid and \PTLNadigRank{} for absolute-effect rank,
with positive 90\% intervals in both source directions. This replication
  supports the ordering-level phenomenon without being used to tune the
  Frangieh decision association. Matched-universe sensitivity for this
  replication is reported in Appendix C, while the complete two-direction
  full-size surface remains in the figure source table.

A second raw-cell replication reconstructs the complete hierarchy in the Tian
iPSC and day-7 neuron CRISPR screens~\citep{tian2019crispri}, rather than using an aggregate rank
audit. We froze 26 shared single-gene perturbations, a 1,536-gene panel chosen
from source controls only, and a 40-cell budget before evaluating either
direction. Five of six direction--metric intervals were positive. From iPSC
to neuron, $D_{\rm adj}$ was $0.032$ [0.005, 0.060], $0.010$
[$-0.010$, 0.030], and $0.016$ [0.004, 0.028] for delta cosine, Systema and
absolute-effect rank, respectively; the reverse direction gave $0.049$
[0.018, 0.079], $0.042$ [0.025, 0.060], and $0.013$ [0.001, 0.025]. The
direction-specific estimates support context-dependent ordering transport;
their asymmetry is reported descriptively rather than treated as a separate
biological conclusion.
The complete six-row ledger and source-frozen construction are in Appendix C.

\subsection{Source-frozen predictor families preserve the qualitative finding}
As one matched predictor-family sensitivity, we also evaluated a source-only
RBF kernel-ridge family on the same Frangieh surface. Its source-only penalty
selection required an expanded grid. The matched transport profile had
Spearman $\rho=0.983$ with the primary bilinear-ridge profile (within-metric
profile correlations $0.829$--$1.000$; mean row-wise prediction-vector cosine
$0.733$ in the Supplement), and adjusted-discrepancy point and interval signs
agreed on all 18 transfer--metric rows. This is descriptive evidence that the
profile is not an immediate artifact of one parameterization, not predictor-
family equivalence or broader generalization.

We then ran the official GEARS graph predictor under the same strict
information contract. For each Frangieh source context, five label-disjoint
folds and all graph construction, fitting and tuning used source cells only;
13 labels absent from the official interaction graph were excluded before
splitting, leaving 230 perturbations and 8,229 evaluation genes. The frozen
out-of-fold vectors were evaluated by the same 30-seed, two-pseudoreplicate
raw-cell bridge. All 18 directed transfer--metric rows retained positive
adjusted-discrepancy intervals. Relative to the bilinear profile, GEARS had
perfect point-sign agreement but lower magnitude concordance (mean
within-metric Spearman $\rho=0.371$: $-0.086$, $1.000$ and $0.200$ for delta
cosine, Systema centroid and absolute-effect rank). Thus a modern graph model
reproduces the existence of transport while showing that its detailed geometry
is predictor- and metric-dependent; it does not establish family equivalence
or superiority. The complete surface is reported in Supplementary Appendix H.

\FloatBarrier
\subsection{Cell budget determines when reordering is resolvable}
We next hold the perturbation universe fixed and resample 10, 20, 40, 80 and
160 cells per perturbation--context pseudoreplicate, together with a matched
full-depth endpoint. The signed $D_{\rm adj}(n)$ trajectory is retained. The
first budget within ten percent of the full-depth value summarizes resolution;
detectability requires a positive lower endpoint of the 90\% interval. These
are metric- and transfer-specific summaries, not a universal cell threshold.
We additionally compare each transport statistic with a matched same-context
pseudo-context null at every depth, because the canonical within-context
measurement term is not itself a same-context $D_{\rm adj}$ estimate.

\begin{figure*}[t]
  \centering
  \includegraphics[width=\textwidth]{../../results/figures/reliability_transportability/reliability_transportability_fig3_measurement_boundary.pdf}
  \caption{\textbf{Cell budget defines a finite-measurement resolution boundary.} (A)
  Signed depth trajectories for the primary Ctrl--IFN$\gamma$ contrast with
  other transfers as pale paths. (B) The six-transfer resolution map; circles
  mark detectability and squares mark the resolution budget, with open glyphs
  denoting not reached. (C) Cross-context transport and matched same-context
  split-half pair-law null margins across depth for the five prespecified
  comparator families. This is a finite-measurement diagnostic, not the
  canonical per-seed $D_{\rm adj}$ summary.
  The full-depth endpoint is the fixed-160 eligibility universe ($n=139$),
  not the wider full-size minimum-20 universe used for Figs.~1--2. Signed
  finite-sample estimates are retained, so null diagnostics can be negative.
  Detectability and quantitative resolution are separate experimental
  requirements under the declared resampling protocol.}
  \label{fig:boundary}
\end{figure*}

The trajectories rise at different rates for the three metrics, and several
source--contrast cells remain unresolved at the available budgets. Increasing
cell budget therefore improves resolution under this protocol without producing a single portable
cutoff. This measurement boundary is a prerequisite for interpreting a
context-associated ordering difference beyond the measured sampling floor. The
matched-null analysis shows positive cross-minus-null margins for every
comparator at every depth, but does not support presenting $D_{\rm adj}$ as a
universally superior rank statistic; its role is the explicit tie-aware,
measurement-adjusted estimand. The depth comparator audit is reported in
the Supplementary finite-measurement audit.
The depth definitions and all transfer-specific results are in
Appendix D.

\FloatBarrier
\subsection{Context-dependent reordering changes shortlist decisions}
To connect the ordering to an experimental action, each source context selects
the lowest-risk perturbations at budgets $k/n\in\{0.05,0.10,0.20,0.50\}$ and
the target context evaluates the selected set. We report target top-$k$
retention and normalized regret under both a random-reference and a bounded
worst-oracle normalization.

\begin{figure*}[t]
  \centering
  \includegraphics[width=\textwidth]{../../results/figures/reliability_transportability/reliability_transportability_fig4_decision_consequence.pdf}
  \caption{\textbf{Context-dependent ordering changes experimental priorities.}
  (A) Detailed example of delta-cosine top-$k$ replacement under the
  Ctrl-to-IFN$\gamma$ shift (6 of 25 shortlist positions are retained; 19 are
  dropped and 19 are new). (B)
  The 18 full-depth transfer--metric observations; marker shape identifies
  metric and pale connectors identify directed transfer. The pointwise error
  bars are canonical 90\% intervals; the descriptive association and its
  unordered-context-pair sensitivity are reported in the text, and no fitted
  line is implied. (C) Unordered-context-pair-held-out incremental prediction
  of normalized regret beyond metric identity; intervals are block-bootstrap
  diagnostics.}
  \label{fig:decision}
\end{figure*}

The primary full-depth link contains six directed source--target surfaces and
three metrics, hence 18 descriptive rows. Across all rows, adjusted reordering
and normalized regret have pooled Spearman $\rho=\PTLDecisionRhoDisplay{}$;
after ranking within metric, the association is
$\rho=\PTLWithinMetricRho{}$.
the three unordered context-pair and source-context sensitivities are reported
separately. The association is a decision-level summary of the declared
surface, not a claim that every predictor or biological system will have the
same slope.

An explicit unordered-context-pair leave-one-out comparison separates metric identity,
mean-fidelity shift, and ordering transport. Adding the matched mean-risk shift
to metric identity reduces OOF MAE by $\PTLIncrementalMeanDeltaMAE{}$
($\PTLIncrementalMeanDeltaMAECI{}$). Adding $D_{\rm adj}$ increased OOF MAE
from 0.087 to 0.092, corresponding to an improvement score of
$\PTLIncrementalDAdjBeyondMeanDeltaMAE{}$
($\PTLIncrementalDAdjBeyondMeanDeltaMAECI{}$); the two nested contrasts are
shown separately in Fig.~\ref{fig:decision}. Thus the pooled association is not
interpreted as a unique causal contribution of one scalar ordering summary.
The corresponding unordered-pair block-bootstrap table is reported in the
Supplement; no headline permutation $p$-value is used for only three blocks.

The shortlist view makes the mechanism concrete: perturbations that occupy the
source top-$k$ need not occupy the target top-$k$, so a source ranking can lose
retention and incur regret even when the predictor itself has not changed. The
budget-wise association summaries show that the link is not confined to one
selection fraction, while its magnitude depends on the transfer and metric.
The four-budget, two-normalization robustness summary is reported in
Appendix D. A metric-matched mean-risk control, using the same eligible labels
and frozen predictions, shows that context-average fidelity and reliability-
order transport are empirically non-equivalent (Supplementary Fig.~S3).

Additional target-informed and prospective source-only analyses are reported in
the Supplement. They are kept separate from the post-hoc transport claim: the
target-informed audit uses target outcomes, whereas the prospective analysis
uses source-observable features only and reserves target outcomes exclusively
for evaluation.

\FloatBarrier
\section{Discussion}
The central result is that perturbation reliability is an environment-dependent
ordering, not a fixed property of a perturbation. With predictions held fixed,
matched biological context shifts reordered realized risk beyond a same-context
sampling floor, and the direction recurs in both directions of an independent
HepG2--Jurkat context pair. The result is therefore an empirical
reliability-transport object: it does not claim that every predictor or
biological system must reorder in the same way.

Unlike response-prediction benchmarks~\citep{lotfollahi2019scgen,
lotfollahi2023cpa,bunne2023cellot,roohani2023gears,cui2024scgpt,
generalizable2025benchmark}, our object asks whether perturbation-level error
ordering remains useful after biology changes. It is also distinct from
experimental perturbation reproducibility~\citep{wang2026reliable}: predictor,
metric and context are part of the estimand.

Cell budget separates detectability from resolution of transport magnitude.
The tie-aware support rule and matched same-context control make this distinction
explicit. In the deterministic tie-free limit, $D_{\rm adj}$ reduces to the
normalized Kendall inversion fraction; here it extends the same question to
uncertain pair states under finite repeated measurement.

Source-ranked top-$k$ choices lose target retention when reliability moves, but
this is a diagnostic rather than a decision-focused learner
~\citep{donti2017task,wilder2019melding,elmachtoub2022smart}. The 18 rows arise
from only three unordered context-pair blocks. Metric-stratified associations
are weaker than the pooled surface; mean-risk shift provides the only modest OOF improvement
($\PTLIncrementalMeanDeltaMAE$; $\PTLIncrementalMeanDeltaMAECI$) and the
additional $D_{\rm adj}$ contrast beyond mean shift does not add held-out value
($\PTLIncrementalDAdjBeyondMeanDeltaMAE$;
$\PTLIncrementalDAdjBeyondMeanDeltaMAECI$). We therefore do not claim that one
scalar universally predicts regret.

\paragraph{Scope and testable extensions.}
The primary decision surface is limited to one bilinear-ridge predictor, three
metrics and three unordered Frangieh context pairs. Matched RBF, neural and
GEARS analyses test predictor sensitivity but are not pooled with it; the Nadig
and Tian reconstructions replicate biological transport but do not enlarge the
decision sample. TxPert lacked a source-matched Frangieh checkpoint, and the
outcome-blind Feng aggregate audit was neutral, so neither is promoted as a
positive replication. The next test is a prospectively frozen, label-disjoint
analysis across additional contexts and modalities. This scope is important as
context diversity becomes central to virtual-cell evaluation
\citep{dibaeinia2026context}.

\paragraph{Reproducibility.}
The \href{https://anonymous.4open.science/r/PTL-B377/}{anonymous core release}
provides fixed splits, source tables and commands for the four main figures,
Supplementary Figs.~S1--S7 and compact tables. Appendices A, B and I specify the
data, estimators and regeneration path; engineering artifacts are excluded from
manuscript evidence.

\paragraph{Code and data availability.}
The same anonymous release contains analysis and figure code, manuscript
sources and compact result products. Raw matrices are not redistributed.
Frangieh and Nadig data are available from the public scPerturb v1.4
\href{https://zenodo.org/records/13350497}{Zenodo record}; the repository
record; Feng data are available from its Cell Genomics record
\citep{feng2026crispri}. The repository README lists exact inputs and commands.
\paragraph{AI use.}
AI-assisted tools supported literature triage, mathematical and experimental
brainstorming, implementation and debugging of analysis code, figure-script
development, interpretation checks, and language editing. The authors designed
the study, inspected all generated analyses, verified citations and numerical
claims against source artifacts, and take responsibility for the manuscript.

\bibliographystyle{../../third_party/iclr2027/iclr2027_conference}
\bibliography{../../manuscript/references}
\end{document}
